%%%%%%%%%%%%%%%%%%%%%%%%%%%%%%%%%%%%%%%%%%%%%%%%%%%%%%%%%%%%
%
% Welcome to overleaf, (formerly writeLaTeX) --- just edit 
% your LaTeX on the left, and we'll compile it for you on 
% the right. If you give someone the link to this page, they 
% can edit at the same time. If you create a free overleaf
% account, all your projects will be saved.  See the help 
% menu above for more info. Enjoy!
% -----------------------------------------------------------
% My thanks to Dana Ernst of Northern Arizona University 
% for sharing his template with me. This is largely his work.
% ----------------------------------------------------------- 
% This is all preamble stuff that you don't have to worry about.
% Head down to where it says "Start here"
% -----------------------------------------------------------
%%%%%%%%%%%%%%%%%%%%%%%%%%%%%%%%%%%%%%%%%%%%%%%%%%%%%%%%%%%%

\documentclass[12pt]{article}
\usepackage[margin=1in]{geometry} 
\usepackage{amsmath,amsthm,amssymb}
\usepackage{graphicx}
\usepackage{hyperref}
 
\newenvironment{theorem}[2][Theorem]{\begin{trivlist}
\item[\hskip \labelsep {\bfseries #1}\hskip \labelsep {\bfseries #2.}]}{\end{trivlist}}
\newenvironment{lemma}[2][Lemma]{\begin{trivlist}
\item[\hskip \labelsep {\bfseries #1}\hskip \labelsep {\bfseries #2.}]}{\end{trivlist}}
\newenvironment{conjecture}[2][Conjecture]{\begin{trivlist}
\item[\hskip \labelsep {\bfseries #1}\hskip \labelsep {\bfseries #2.}]}{\end{trivlist}}
\newenvironment{question}[2][Question]{\begin{trivlist}
\item[\hskip \labelsep {\bfseries #1}\hskip \labelsep {\bfseries #2.}]}{\end{trivlist}}
\newenvironment{corollary}[2][Corollary]{\begin{trivlist}
\item[\hskip \labelsep {\bfseries #1}\hskip \labelsep {\bfseries #2.}]}{\end{trivlist}}
\newenvironment{definition}[2][Definition]{\begin{trivlist}
\item[\hskip \labelsep {\bfseries #1}\hskip \labelsep {\bfseries #2.}]}{\end{trivlist}}
\newenvironment{problem}[2][Problem]{\begin{trivlist}
\item[\hskip \labelsep {\bfseries #1}\hskip \labelsep {\bfseries #2.}]}{\end{trivlist}}

\begin{document}
 
% --------------------------------------------------------------
%                         Start here
% --------------------------------------------------------------
 
\title{Inquiry \& Proof} 
\author{Your Name} % replace with your name
%\date{A fixed date} % for a fixed date, uncomment this line and replace date, otherwise, today's date will be used.
\maketitle

\noindent Information about collaboration should go here.  If you feel it is appropriate to include any other comments as an introduction, put them here too.

\begin{problem}{x.yz} %You can use theorem, exercise, problem, or question here.  Modify x.yz to be whatever number you are explaining/proving.
Delete this text and write your statement here. Replace x.yz with the class's numbering convention for that item.
\end{problem}
 
\begin{proof}  
Blah, blah, blah. Delete this and write your proof here.
\end{proof}
 
\bigskip

\noindent You'll generally be turning in one item per page, but here is another example, just so we can showcase a few more features of \LaTeX.   

\bigskip

\begin{theorem}{w.xy}
This is a theorem about flying pigs.
\end{theorem}
 
\begin{proof}
Blah, blah, blah.  This is another proof for you to delete. It has some math symbols in it:  Let $\alpha \in (A \cup B) \cap \{\text{flying pigs}\}$ and $\mathbb{Z} \subset \mathbb{R}$.  

\begin{description}
\item[Case 1:] Suppose $x=\sin\theta$ ... thus pigs fly.
\item[Case 2:] Suppose $x \neq \sin\theta$ ... therefore, pigs fly.
\end{description}
Since in every case pigs fly, we have proven that pigs do fly.

\end{proof}

\noindent If you want to include a figure, you have to upload the image file to overleaf (or store it in the same folder as your .tex file, if you're not using overleaf).  Look in the .tex source file for the syntax, which has been commented out so it doesn't appear here in the pdf. 

%\begin{figure}[ht]
%\centering
%\includegraphics[width=.25\textwidth]{myimage.png}
%\caption{This is a picture of math.}
%\end{figure}

It is possible to change the way your figure is placed and sized. Play around with the parameters to see what happens.

By the way, when you make figures with GeoGebra, you can easily ``export" as a .png file. Then upload into overleaf.com using the ``Project'' button, then "Add files..." 
  
% --------------------------------------------------------------
%     Leave this at the bottom and replace URL with your own
% --------------------------------------------------------------

\vfill
\small{Read-only document can be found online at \url{https://www.overleaf.com/lotasymbols}} %copy and paste the READ-ONLY URL of your document into this line.  Then others with this pdf can also access the .tex file.  Avoid editing or deleting the project once you've shared this URL. (create a copy of this project to start a revised version of the same.)  I'll use this mostly when a result gets published to the online class journal.
 

% --------------------------------------------------------------
%     You don't have to mess with anything below this line.
% --------------------------------------------------------------
 
\end{document}